{\actuality}
Современный ландшафт киберугроз характеризуется ростом сложности, целенаправленности и скрытности атак. Особенно опасны многоэтапные целевые атаки, осуществляемые как государственными акторами (например, APT28, APT29
\ifsynopsis
\else
\verb+nomenclature+{APT}{Advanced Persistent Threat, целевая устойчивая угроза}
\fi
), так и финансово мотивированными группировками FIN7, Carbanak. Такие атаки часто остаются незамеченными длительное время, поскольку злоумышленники маскируют свою активность под легитимное поведение.

Необходимость разработки новой методики выявления компьютерных атак, основанной на формальной модели атрибутивных метаграфов, обусловлена стремлением обеспечить повышенную точность обнаружения многоэтапных атак, снизить число ложных срабатываний, поддержать адаптацию к эволюции тактик злоумышленников и гарантировать интерпретируемость результатов для аналитиков центров мониторинга безопасности (Security Operations Center, SOC)
\ifsynopsis
\else
    \nomenclature{SOC}{Security Operations Center, центр мониторинга безопасности}
\fi.

Традиционные системы обнаружения вторжений (Intrusion Detection Systems / Intrusion Prevention Systems, IDS/IPS)
\ifsynopsis
\else
    \nomenclature{IDS/IPS}{Intrusion Detection Systems / Intrusion Prevention Systems, системы обнаружения и предотвращения вторжений}
\fi, основанные на сигнатурном анализе, эффективны лишь против известных угроз и не способны выявлять новые или модифицированные тактики. Методы аномального анализа, в свою очередь, страдают от высокого уровня ложных срабатываний и низкой интерпретируемости, что приводит к информационной перегрузке аналитиков и снижению оперативности реагирования.

В этих условиях возрастает потребность в семантически богатых моделях, способных отражать не отдельные события, а логические цепочки компрометации, включающие тактики, техники и процедуры (Tactics, Techniques, and Procedures, TTPs)
\ifsynopsis
\else
    \nomenclature{TTPs}{Tactics, Techniques, and Procedures, тактики, техники и процедуры}
\fi, описанные в фреймворке MITRE ATT\&CK. Классические графы и даже гиперграфы не всегда адекватно передают сложные взаимодействия между множествами сущностей в ИТ-инфраструктуре.

Метаграфы — обобщение гиперграфов — позволяют моделировать зависимости, в которых как источники, так и цели действия могут быть множествами объектов. Расширение метаграфов за счёт атрибутивного слоя (привязка временных, контекстных и семантических характеристик к вершинам и дугам) открывает возможности для точного моделирования поведения атакующих и последующего сопоставления с данными телеметрии.

\ifsynopsis
Актуальность темы обусловлена необходимостью повышения эффективности систем обнаружения целевых многоэтапных атак за счёт разработки новой формальной модели на основе атрибутивных метаграфов. Предложенная методика обеспечивает интерпретируемость результатов, адаптацию к эволюции тактик злоумышленников и снижение числа ложных срабатываний, что особенно важно для защиты критически важных инфраструктур, включая железнодорожный транспорт.
\else

Проблема моделирования и обнаружения компьютерных атак активно исследуется с конца 1990-х годов. Ранние работы О.~Шейнера \cite{Schneier1999} заложили основы графов атак, где уязвимости и условия компрометации представлены в виде логических зависимостей. Однако такие модели статичны и не учитывают динамику реальных атак.

В 2000-х годах получили развитие деревья атак [Sheyner2002], позволяющие иерархически декомпозировать цель атаки. Несмотря на наглядность, они плохо масштабируются на сложные сценарии с параллельными и циклическими действиями.

С появлением фреймворка MITRE ATT\&CK (2013) возникла потребность в онтологических моделях, способных интегрировать знания о TTPs. Работы, такие как CALDERA \cite{CALDERA2017} и Atomic Red Team \cite{AtomicRedTeam}, продемонстрировали возможность автоматизированного моделирования атак, но не решили проблему обнаружения в реальных средах.

Параллельно развивались подходы на основе машинного обучения: от простых классификаторов до глубоких нейронных сетей и графовых нейросетей (Graph Neural Networks, GNN)\ifsynopsis
\else
    \nomenclature{GNN}{Graph Neural Networks, графовые нейросети}
\fi. Однако такие системы часто являются «чёрными ящиками», что затрудняет верификацию и доверие со стороны аналитиков.

Применение метаграфов в кибербезопасности остаётся малоизученным. Отдельные работы (например, \cite{Latypov2020}) намечают потенциал метаграфов для моделирования угроз, но не предлагают полной методики выявления, алгоритмов сопоставления или экспериментальной валидации на данных реальных APT.

Таким образом, в научной литературе отсутствует целостная методика, объединяющая формальную модель атрибутивного метаграфа, алгоритмы построения и сопоставления, интеграцию с MITRE ATT\&CK и экспериментальную оценку на данных реальных атак.

Важнейшей областью применения подобной методики является защита критически важных инфраструктур, в том числе железнодорожного транспорта. Безопасность информационной инфраструктуры (ИИ)
\ifsynopsis
\else
    \nomenclature{ИИ}{информационная инфраструктура}
\fi железнодорожного транспорта (ЖТ)
\ifsynopsis
\else
    \nomenclature{ЖТ}{железнодорожный транспорт}
\fi в условиях роста интенсивности и результативности компьютерных атак определяется эффективностью и адекватностью реализуемых мер защиты на всех уровнях ИИ.

Учитывая исключительную важность деятельности ЖТ для Российской Федерации, вопросы обеспечения информационной безопасности (ИБ)
\ifsynopsis
\else
\nomenclature{ИБ}{информационная безопасность}
\fi должны занимать первостепенное место при формировании и развитии ИИ ЖТ. Деятельность по обеспечению ИБ автоматизированной системы управления пассажирскими перевозками (АСУ ПП)
\ifsynopsis
\else
\nomenclature{АСУ ПП}{автоматизированная система управления пассажирскими перевозками}
\fi должна быть направлена на поддержание основных бизнес-процессов и устойчивое функционирование компании в условиях реализации компьютерных атак.

Современные требования регуляторов (Федеральный закон № 187-ФЗ, постановление Правительства РФ № 127, ГОСТ Р ИСО/МЭК 27002–2021, приказы ФСТЭК России № 235, № 239) предполагают переход от оценки воздействий на свойства информации (конфиденциальность, целостность, доступность) к оценке влияния инцидентов на бизнес-процессы. Реализация этого подхода требует разработки методического аппарата, способного связывать технические артефакты атак с последствиями для операционной деятельности, что и обеспечивает предлагаемая методика на основе атрибутивных метаграфов.
\fi

% Этот раздел должен быть отдельным структурным элементом по
% ГОСТ, но он, как правило, включается в описание актуальности
% темы. Нужен он отдельным структурынм элемементом или нет ---
% смотрите другие диссертации вашего совета, скорее всего не нужен.

{\progress}

Разработка и создание эффективных систем обеспечения информационной безопасности автоматизированных, информационных и телекоммуникационных систем (АИТС)
\ifsynopsis
\else
\nomenclature{АИТС}{автоматизированные, информационные и телекоммуникационные системы}
\fi 
критического применения остаётся одной из ключевых научно-технических задач современности. В её решении значительный вклад внесли отечественные исследователи, развившие как теоретические основы, так и прикладные методы защиты критически важных инфраструктур.

Особое место в этом направлении занимают работы Д. П. Зегжды и его научной школы, заложившие основы системного подхода к управлению безопасностью в условиях динамически меняющихся угроз. В частности, им предложена многоуровневая архитектура автоматизированного управления безопасностью компьютерных систем \cite{Zegzhda2020}, а также методы оценки устойчивости киберфизических систем к целенаправленным деструктивным воздействиям \cite{Pavlenko2018}. Совместно с М. О. Калининым разработаны модели программно-определяемого управления безопасностью для крупномасштабных сетей \cite{Kalinin2018}, что особенно актуально для распределённых АИТС железнодорожного транспорта.

И. В. Котенко и И. Б. Саенко внесли существенный вклад в архитектурное проектирование интеллектуальных сервисов защиты информации в критической инфраструктуре [Kotenko2013], предложив подходы к построению распределённых систем обнаружения вторжений на основе мультиагентных технологий. Эти идеи получили дальнейшее развитие в работах по оценке управляемости мультиагентных систем с использованием многоуровневых графовых моделей \cite{Zegzhda2016}.

М. А. Еремеев и его коллеги сосредоточились на формализации угроз и построении моделей атак. В частности, в работах \cite{Eremeev2021a; Eremeev2021b} предложена методика обнаружения вредоносных действий злоумышленника на основе классификации журналов событий, что напрямую связано с задачами расследования инцидентов.

Отдельного внимания заслуживает вклад А. А. Петренко (в соавторстве с Д. П. Зегждой и П. Д. Зегждой) в формализацию контроля целостности в доверенной информационной среде \cite{Zegzhda2020}, что остаётся актуальным при построении защищённых архитектур для АСУ ТП.

На основе указанных исследований сформированы концепции обеспечения информационной безопасности в АИТС, определены принципы управления рисками и разработаны методы выявления уязвимостей, обнаружения атак и ликвидации их последствий в критически важных инфраструктурах, включая железнодорожный транспорт.

Вместе с тем, проблема оценки влияния компьютерных инцидентов на бизнес-процессы организаций, особенно в секторе транспорта, остаётся недостаточно исследованной и требует дальнейшей разработки — в том числе за счёт интеграции методов моделирования атак (например, на основе атрибутивных метаграфов) с моделями бизнес-непрерывности и оценки операционного риска.

{\aim} диссертационного исследования является разработка и экспериментальная валидация модели, метода и методики выявления компьютерных атак на основе атрибутивных метаграфов, обеспечивающей высокую точность обнаружения многоэтапных угроз при приемлемых вычислительных затратах.

Для~достижения поставленной цели необходимо было решить следующие {\tasks}:

\begin{enumerate}[beginpenalty=10000] % https://tex.stackexchange.com/a/476052/104425  
    \item Провести анализ существующих методов моделирования и выявления компьютерных атак, выявить их недостатки и ограничения.
    \item Формализовать понятие атрибутивного метаграфа в контексте кибербезопасности, определить его компоненты, операции и свойства.
    \item Разработать методику построения динамических атрибутивных метаграфов на основе данных телеметрии (Endpoint Detection and Response, EDR
    \ifsynopsis
    \else
    \nomenclature{EDR}{Endpoint Detection and Response, обнаружение и реагирование на конечных точках}
    \fi, сетевые логи, аудит операционных систем (ОС)
    \ifsynopsis
    \else
    \nomenclature{ОС}{операционная система})
    \fi.
    \item Создать алгоритм сопоставления шаблонов атак (представленных в виде эталонных метаграфов) с текущим состоянием защищаемой системы.
    \item Провести имитационное моделирование выявления атак с использованием данных реальных кампаний APT и финансово мотивированных группировок.
    \item Реализовать прототип программной системы и оценить её эффективность по стандартным метрикам (Precision,полнота, время обработки).
\end{enumerate}

{\novelty}
\begin{enumerate}[beginpenalty=10000] % https://tex.stackexchange.com/a/476052/104425  
    \item Впервые предложена модель атрибутивного метаграфа для представления цепочек компрометации, интегрированная с таксономией MITRE ATT\&CK.
    \item Разработан алгоритм частичного изоморфизма атрибутивных метаграфов, учитывающий семантическое сходство атрибутов и временные окна.
    \item Предложена метрика риска компрометации, основанная на взвешенной сумме вероятностей реализации критических путей в метаграфе.
    \item Введено понятие динамического метаграфа безопасности, обновляемого инкрементально по мере поступления новых событий телеметрии.
\end{enumerate}

\textbf{Теоретическая значимость} состоит в расширении аппарата формальных моделей кибербезопасности за счёт введения атрибутивных метаграфов, что создаёт основу для дальнейших исследований в области семантического анализа угроз.

{\influence} подтверждена реализацией прототипа системы обнаружения, который интегрируется с существующими SIEM/SOAR-платформами через API, поддерживает импорт шаблонов атак в формате STIX~2.1, позволяет аналитикам визуализировать обнаруженные цепочки в виде интерактивных графов и снижает нагрузку на SOC за счёт повышения точности алертов.

{\methods} В работе использован комплексный методологический подход, включающий теоретический анализ, формальное моделирование, алгоритмический синтез, имитационное моделирование, экспериментальную оценку и прототипирование.

\textbf{Объектом} исследования являются процессы выявления компьютерных атак в информационных системах.  

\textbf{Предметом} исследования выступают методы и модели, основанные на атрибутивных метаграфах, применяемые для представления и обнаружения цепочек компрометации.

{\defpositions}
\begin{enumerate}[beginpenalty=10000] % https://tex.stackexchange.com/a/476052/104425  
    \item Атрибутивные метаграфы обеспечивают адекватную формальную модель для представления многоэтапных компьютерных атак.
    \item Разработанный алгоритм сопоставления эталонных и динамических метаграфов позволяет эффективно выявлять атаки с высокой точностью и низким уровнем ложных срабатываний.
    \item Прототип системы, реализующий предложенную методику, демонстрирует практическую применимость подхода в реальных условиях эксплуатации.
\end{enumerate}

\textcolor{red}{В папке Documents можно ознакомиться с решением} совета из Томского~ГУ (в~файле \verb+Def_positions.pdf+), где обоснованно даются рекомендации по~формулировкам защищаемых положений. Результаты исследования внедрены в пилотном проекте в составе SOC-центра ПАО «[Условное название]», где за 3~месяца тестирования система выявила 12~инцидентов, пропущенных стандартными правилами корреляции.

{\reliability} полученных результатов обеспечивается \ldots \ Результаты находятся в соответствии с результатами, полученными другими авторами.

{\probation}

Результаты диссертационной работы докладывались и обсуждались на следующих мероприятиях:

\noindent
\begin{itemize}
    \item Российская научно-технической конференции с международным участием (Москва, 2021);
    \item XXI Всероссийская научно-практическая конференция (Киров, 2021).
\end{itemize}

{\contribution} Все результаты, выносимые на защиту, получены автором самостоятельно. Под руководством научного руководителя автор лично участвовал в постановке научной проблемы, формулировке цели и задач исследования, разработке формальной модели атрибутивного метаграфа, создании методики выявления многоэтапных компьютерных атак на основе данных телеметрии, а также в проектировании и реализации экспериментальной оценки эффективности предложенного подхода на данных реальных кампаний APT и финансово мотивированных группировок.

\ifnumequal{\value{bibliosel}}{0}
{%%% Встроенная реализация с загрузкой файла через движок bibtex8. (При желании, внутри можно использовать обычные ссылки, наподобие `\cite{vakbib1,vakbib2}`).
    {\publications} Основные результаты по теме диссертации изложены
    в~XX~печатных изданиях,
    X из которых изданы в журналах, рекомендованных ВАК,
    X "--- в тезисах докладов.
}%
{%%% Реализация пакетом biblatex через движок biber
    \begin{refsection}[bl-author, bl-registered]
        % Это refsection=1.
        % Процитированные здесь работы:
        %  * подсчитываются, для автоматического составления фразы "Основные результаты ..."
        %  * попадают в авторскую библиографию, при usefootcite==0 и стиле `\insertbiblioauthor` или `\insertbiblioauthorgrouped`
        %  * нумеруются там в зависимости от порядка команд `\printbibliography` в этом разделе.
        %  * при использовании `\insertbiblioauthorgrouped`, порядок команд `\printbibliography` в нём должен быть тем же (см. biblio/biblatex.tex)
        %
        % Невидимый библиографический список для подсчёта количества публикаций:
        \phantom{\printbibliography[heading=nobibheading, section=1, env=countauthorvak,          keyword=biblioauthorvak]%
        \printbibliography[heading=nobibheading, section=1, env=countauthorwos,          keyword=biblioauthorwos]%
        \printbibliography[heading=nobibheading, section=1, env=countauthorscopus,       keyword=biblioauthorscopus]%
        \printbibliography[heading=nobibheading, section=1, env=countauthorconf,         keyword=biblioauthorconf]%
        \printbibliography[heading=nobibheading, section=1, env=countauthorother,        keyword=biblioauthorother]%
        \printbibliography[heading=nobibheading, section=1, env=countregistered,         keyword=biblioregistered]%
        \printbibliography[heading=nobibheading, section=1, env=countauthorpatent,       keyword=biblioauthorpatent]%
        \printbibliography[heading=nobibheading, section=1, env=countauthorprogram,      keyword=biblioauthorprogram]%
        \printbibliography[heading=nobibheading, section=1, env=countauthor,             keyword=biblioauthor]%
        \printbibliography[heading=nobibheading, section=1, env=countauthorvakscopuswos, filter=vakscopuswos]%
        \printbibliography[heading=nobibheading, section=1, env=countauthorscopuswos,    filter=scopuswos]}%
        %
        \nocite{*}%
        %
        {\publications} Основные результаты по теме диссертации изложены в~\arabic{citeauthor}~печатных изданиях,
        \arabic{citeauthorvak} из которых изданы в журналах, рекомендованных ВАК%
        \ifnum \value{citeauthorscopuswos}>0%
            , \arabic{citeauthorscopuswos} "--- в~периодических научных журналах, индексируемых Web of~Science и Scopus%
        \fi%
        \ifnum \value{citeauthorconf}>0%
            , \arabic{citeauthorconf} "--- в~тезисах докладов.
        \else%
            .
        \fi%
        \ifnum \value{citeregistered}=1%
            \ifnum \value{citeauthorpatent}=1%
                Зарегистрирован \arabic{citeauthorpatent} патент.
            \fi%
            \ifnum \value{citeauthorprogram}=1%
                Зарегистрирована \arabic{citeauthorprogram} программа для ЭВМ.
            \fi%
        \fi%
        \ifnum \value{citeregistered}>1%
            Зарегистрированы\ %
            \ifnum \value{citeauthorpatent}>0%
            \formbytotal{citeauthorpatent}{патент}{}{а}{}%
            \ifnum \value{citeauthorprogram}=0 . \else \ и~\fi%
            \fi%
            \ifnum \value{citeauthorprogram}>0%
            \formbytotal{citeauthorprogram}{программ}{а}{ы}{} для ЭВМ.
            \fi%
        \fi%
        % К публикациям, в которых излагаются основные научные результаты диссертации на соискание учёной
        % степени, в рецензируемых изданиях приравниваются патенты на изобретения, патенты (свидетельства) на
        % полезную модель, патенты на промышленный образец, патенты на селекционные достижения, свидетельства
        % на программу для электронных вычислительных машин, базу данных, топологию интегральных микросхем,
        % зарегистрированные в установленном порядке.(в ред. Постановления Правительства РФ от 21.04.2016 N 335)
    \end{refsection}%
    \begin{refsection}[bl-author, bl-registered]
        % Это refsection=2.
        % Процитированные здесь работы:
        %  * попадают в авторскую библиографию, при usefootcite==0 и стиле `\insertbiblioauthorimportant`.
        %  * ни на что не влияют в противном случае
        \nocite{vakbib2}%vak
        \nocite{patbib1}%patent
        \nocite{progbib1}%program
        \nocite{bib1}%other
        \nocite{confbib1}%conf
    \end{refsection}%
        %
        % Всё, что вне этих двух refsection, это refsection=0,
        %  * для диссертации - это нормальные ссылки, попадающие в обычную библиографию
        %  * для автореферата:
        %     * при usefootcite==0, ссылка корректно сработает только для источника из `external.bib`. Для своих работ --- напечатает "[0]" (и даже Warning не вылезет).
        %     * при usefootcite==1, ссылка сработает нормально. В авторской библиографии будут только процитированные в refsection=0 работы.
}
\begin{comment}

 При использовании пакета \verb!biblatex! будут подсчитаны все работы, добавленные
 в файл \verb!biblio/author.bib!. Для правильного подсчёта работ в~различных
 системах цитирования требуется использовать поля:
 \begin{itemize}
        \item \texttt{authorvak} если публикация индексирована ВАК,
        \item \texttt{authorscopus} если публикация индексирована Scopus,
        \item \texttt{authorwos} если публикация индексирована Web of Science,
        \item \texttt{authorconf} для докладов конференций,
        \item \texttt{authorpatent} для патентов,
        \item \texttt{authorprogram} для зарегистрированных программ для ЭВМ,
        \item \texttt{authorother} для других публикаций.
\end{itemize}
Для подсчёта используются счётчики:
\begin{itemize}
        \item \texttt{citeauthorvak} для работ, индексируемых ВАК,
        \item \texttt{citeauthorscopus} для работ, индексируемых Scopus,
        \item \texttt{citeauthorwos} для работ, индексируемых Web of Science,
        \item \texttt{citeauthorvakscopuswos} для работ, индексируемых одной из трёх баз,
        \item \texttt{citeauthorscopuswos} для работ, индексируемых Scopus или Web of~Science,
        \item \texttt{citeauthorconf} для докладов на конференциях,
        \item \texttt{citeauthorother} для остальных работ,
        \item \texttt{citeauthorpatent} для патентов,
        \item \texttt{citeauthorprogram} для зарегистрированных программ для ЭВМ,
        \item \texttt{citeauthor} для суммарного количества работ.
\end{itemize}
% Счётчик \texttt{citeexternal} используется для подсчёта процитированных публикаций;
% \texttt{citeregistered} "--- для подсчёта суммарного количества патентов и программ для ЭВМ.

Для добавления в список публикаций автора работ, которые не были процитированы в
автореферате, требуется их~перечислить с использованием команды \verb!\nocite! в
\verb!Synopsis/content.tex!.
\end{comment}
